% META: {"id":"1SPE-DERLOCAL-CO-040","chapitre":"1SPE-DERIVATION-LOCAL","type_objet":"corrige","exercice_id":"1SPE-DERLOCAL-EX-040","status":"approved","origin":"generated","status_history":["generated","audited","accepted","approved"],"release_acceptance":"RELEASE_OWNER_BATCH_ACCEPTANCE","acceptance_closure_digest":"sha256:9b3ccf9a81c5520fbb7e03b7b2d3e7bf2057a02834b6d908bf3be5f49f3b553f","acceptance_receipt_digest":"sha256:4fad86f0282e0fa4e0419cca1ad6379ae7af5a280c04a4a90880f90a1c044242"}
% BEGIN-VERIFY
% from sympy import symbols, Rational
% t = symbols('t')
% T = 100 - 2*t**2
% assert T.subs(t, 0) == 100
% assert T.subs(t, 5) == 50
% taux = Rational(T.subs(t, 5) - T.subs(t, 0), 5 - 0)
% assert taux == -10
% END-VERIFY
\begin{corrige}{1SPE-DERLOCAL-EX-040}
\textbf{1.} On calcule les températures aux instants $t = 0$ et $t = 5$ :
\[
T(0) = 100 - 2 \times 0^2 = 100
\]
\[
T(5) = 100 - 2 \times 5^2 = 100 - 50 = 50
\]

\textbf{2.} Le taux de variation de $T$ entre $t = 0$ et $t = 5$ est :
\[
\tau = \frac{T(5) - T(0)}{5 - 0} = \frac{50 - 100}{5} = \frac{-50}{5} = -10
\]

\textbf{3.} Le taux de variation est négatif ($\tau = -10$), ce qui signifie que la température du four diminue après extinction. En moyenne, la température baisse de $\boxed{10}$ degrés Celsius par minute entre les instants $t = 0$ et $t = 5$. Le signe négatif traduit le refroidissement du four.
\end{corrige}
